\begin{lema}[Escalado después de una jugada]\label{lem:hijo-fisico}
Sea un padre regular de tamaño \(m\) cuyo punto físico \(p^{Q}\) es positivo en su soporte
binario, tiene márgenes de fila y columna uno, \(p^{Q}_{ic}\le C/m\) y brecha de incidencia
al menos \(\gamma>0\). Fijada una casilla legal \(e\), fórmese el menor de torres \(M_e\)
y prohíbase cualquier subconjunto de las demás casillas de las dos diagonales de \(e\).
Para \(m\) suficientemente grande, todas esas máscaras de tamaño \(m-1\) tienen escalado en
una misma clase de difusión y brecha, con constantes que sólo dependen de \(C,\gamma\).
En particular esto vale para el hijo completo \(M_e^{\rm fis}\) en la primera salida.
\end{lema}
\begin{proof}
Eliminamos una columna en cada fila retenida y hasta una casilla por cada diagonal.
La pérdida en cada fila es a lo sumo \(3C/m\), y lo mismo vale por columna. Sea \(R\) la
restricción de \(p^Q\) después de las eliminaciones y sea \(\mathcal W\) el lado con menos vértices
de un corte del grafo bipartito restante. En el padre, Rayleigh sobre el indicador centrado
da capacidad al menos \(\gamma |\mathcal W|/2\), pues \(|\mathcal W|\le m\).
Queda en \(R\) una capacidad de al menos \((\gamma/2-3C/m)|\mathcal W|\).
Oriente las aristas de filas a columnas. Los déficits de fila y los déficits de columna
con signo opuesto tienen suma cero y cada uno tiene módulo a lo sumo \(3C/m\).
Con \(\lambda=(3C/m)/(\gamma/2-3C/m)<1\), cada corte soporta la demanda neta
con capacidades simétricas \(\lambda R\). El criterio de flujo factible~\cite{gale} da una corrección
firmada
y una matriz bistocástica \((1-\lambda)R\le p'\le(1+\lambda)R\).
Su conductancia es al menos \((1-\lambda)(\gamma/2-3C/m)\). La desigualdad de Cheeger~\cite{lsok} da una
brecha
positiva uniforme para el laplaciano normalizado.

Falta transferir estas cotas de una tabla factible al escalado del soporte binario.
En tamaño \(k\), sea \(a\) bistocástica, positiva en el soporte, con \(a_{ic}\le C'/k\)
y brecha \(\gamma'\). El máximo de entropía en ese soporte es interior por la factibilidad
estricta. Sus ecuaciones dan el escalado \(D\) con \(\log(kD_{ic})=u_i-u_c\).
Por los márgenes, \(\sum a_{ic}\log(kD_{ic})=\sum D_{ic}\log(kD_{ic})\ge0\).
La parte positiva de la primera suma es a lo sumo \(\sum a_{ic}kD_{ic}\le C'k\).
Al restar una mediana de los potenciales e integrar los cortes de nivel se obtiene
\[
 \frac{\gamma'}2\sum_z|u_z|\le\sum_{ic}a_{ic}|u_i-u_c|\le2C'k .
\]
Fije \(T>8(C')^2/\gamma'\). Cada línea tiene al menos \(k/C'\) vecinos, y en todo el grafo
menos de \(k/(2C')\) vértices tienen \(|u_z|>T\). La ecuación de fila
\(1=e^{u_i}\sum_{c\sim i} e^{-u_c}/k\) da \(e^{u_i}\le2C'e^T\), y la de columna da
\(e^{-u_c}\le2C'e^T\). Las mismas ecuaciones con las cotas superiores dan cotas inferiores
recíprocas. Así \(d_-/k\le D_{ic}\le d_+/k\) en el soporte y
\(D\ge(d_-/C')a\) conserva una brecha positiva por Rayleigh.
Se aplica esto a \(p'\), con \(k=m-1\). Para la máscara del padre se aplica directamente a
\(p^Q\), también en la llegada larga regular. Los márgenes del predictor~\cite{pc} se obtienen
hasta la parada por el Corolario~4.12 y la clase interior de~\cite[def.~2.2 y §§4--5]{tc}.
Si la brecha se enuncia con pesos de respuesta, su comparabilidad con \(p^Q\) en la rama
de vacancias positivas da la hipótesis de Rayleigh.
\end{proof}

El coeficiente \(\coefT=\mathfrak F-1/12\) de (T) depende de las masas marcadas de fila y
columna \(s,t\) y del peso del cruce \(z\). Las expresiones integradas del
apéndice~\ref{app:rook} son holomorfas y acotadas en discos con radios de clase.
En \(s=0\) o \(t=0\) los términos visibles son independientes de \(z\).
Al derivar en \(z\), ambas caras se anulan. Si \(M_0\) acota \(\coefT\) en un
polidisco de radio \(r\) de clase en torno al punto de borrado en \(z\), Cauchy en
\(z\) acota \(|\partial_z\coefT|\) por \(2M_0/r\) en el disco reducido. Schwarz en las
masas \(s,t\) introduce los factores \(|s|/r,|t|/r\). Por tanto
\(\partial_z\coefT=st\,O(1)\) uniformemente, también al integrar en \(z\) a lo largo
del segmento de borrado. Las dos masas valen \(1/k\); en tamaño \(k\) esto da
\(O(k^{-3})\) para el cambio de \(\coefT/k\) por casilla vetada. Retirar una fila y una columna
cambia las masas en \(O(k^{-1})\) y el tamaño inverso en \(O(k^{-2})\), con coste
\(O(k^{-2})\) en \(\coefT/k\). Hay a lo sumo \(2m\) casillas vetadas. El lema proporciona
los mismos márgenes en todos los pasos. Sumando los cambios y usando el resto \(O(m^{-2})\)
de (T) en los dos extremos, obtenemos para todos los hijos del padre regular
\begin{equation}\label{eq:incremento-hijo-fisico}
 |g_T(M_e^{\rm fis})-g_T(M)|\le C/m^2,\qquad |g_T(M_e^{\rm fis})|\le C/(m-1).
\end{equation}
